\section{Proposed Method: {\graphsaint}}

%3.5 pages (End of page 5.5)

Graph sampling based method is motivated by the challenges in scalability (in terms of model depth and graph size). We analyze the bias (Section \ref{sec: bias}) and variance (Section \ref{sec: var}) introduced by graph sampling, and thus, propose feasible sampling algorithms (Section \ref{sec: sampler}). 
We show the applicability of {\graphsaint} to other architectures, both conceptually (Section \ref{sec: discussion}) and experimentally (Section \ref{sec: exp extension}).

In the following, we define the problem of interest and the corresponding notations. 
A GCN learns representation of an un-directed, attributed graph $\G\paren{\V,\E}$, where each node $v\in\V$ has a length-$f$ attribute $\bm{x}_v$.
Let $\bm{A}$ be the adjacency matrix and $\Anorm$ be the normalized one (i.e., $\Anorm=\bm{D}^{-1}\bm{A}$, and $\bm{D}$ is the diagonal degree matrix). 
Let the dimension of layer-$\ell$ input activation be $f^{\paren{\ell}}$. The activation of node $v$ is $\bm{x}_v^{\paren{\ell}}\in\mathbb{R}^{f^{\paren{\ell}}}$, and the weight matrix is $\W[\ell]\in\mathbb{R}^{f^{\paren{\ell}}\times f^{\paren{\ell+1}}}$. Note that $\bm{x}_v=\bm{x}_v^{\paren{1}}$. 
Propagation rule of a layer is defined as follows:

\begin{equation}
\label{eq: vanilla gcn}
    \bm{x}_v^{\paren{\ell+1}} = \sigma\paren{\sum\limits_{u\in\V}\Anorm_{v,u}\paren{\W[\ell]}^{\trans} \bm{x}_u^{\paren{\ell}}}
\end{equation}

where $\Anorm_{v,u}$ is a scalar, taking an element of $\Anorm$. And $\sigma\paren{\cdot}$ is the activation function (e.g., ReLU). 

We use subscript ``s'' to denote parameterd of the sampled graph (e.g., $\G[s]$, $\V[s]$, $\E[s]$).

GCNs can be applied under inductive and transductive settings. While {\graphsaint} is applicable to both, in this paper, we focus on inductive learning. 
It has been shown that inductive learning is especially challenging \citep{graphsage} --- during training, neither attributes nor connections of the test nodes are present. Thus, an inductive model has to generalize to completely unseen graphs. 
%In addition, both the analysis and experiments are under the supervised learning setting. 
% ---- move this to exp as the setup ---- We focus on supervised learning of GCNs. Each training graph node is provided with a ground-truth label. Prediction of the labels is made by an MLP whose inputs are the node embeddings generated by the GCN. 

%Below we define the notations. Let $\G\paren{\V,\E,\X}$ be the training graph with attribute $\X\in\mathbb{R}^{\size{\V}\times f^{\paren{0}}}$. 
%Let $\bm{A}$ be the adjacency matrix of $\G$, and $\Anorm$ be the normalized adjacency matrix ().%(See \cite{gcn,gcn_nips16} for various normalization methods). 
%Let $\xvl{v}{\ell}\in \mathbb{R}^{f^{\paren{\ell}}}$ denote the activation of a node $v$ in GCN layer $\ell$. Let $\X[\paren{\ell}]\in \mathbb{R}^{\size{\V}\times f^{\paren{\ell}}}$ 
%be the activation matrix consisting of $\xvl{v}{\ell}$. Regarding the input layer as layer $0$, we have $\X[\paren{0}]=\X$. 
%Further, let $\bm{y}_v$ be the predicted label of node $v$, and $\overline{\bm{y}}_v$ be the ground truth. Let $\bm{Y}$ and $\overline{\bm{Y}}$ be the label matrix constructed from the corresponding label vectors. Lastly, let the set of weight matrices associated with layer $\ell$ be $\set{\bm{W}^{\paren{\ell}}}$. 
%Unless otherwise specified, symbols with subscript $s$ denote parameters corresponding to a sampled subgraph of $\G$ (e.g., $\G[s]$ is a sampled graph from $\G$, and $\Anorm[s]$ is the normalized adjacency matrix of $\G[s]$). 
%Parameters within parentheses ``$(\cdot)$'' in superscript denote GCN layer index, and those without parentheses denote power of a matrix. 


\subsection{Minibatch by Graph Sampling\label{sec: gsaint}}



\begin{figure}[!htb]
    \centering
    % blue
% \definecolor{c0}{HTML}{e1f5fe}
% \definecolor{c1}{HTML}{b3e5fc}
% \definecolor{c2}{HTML}{4fc3f7}
% \definecolor{c3}{HTML}{4fc3f7}
% \definecolor{c4}{HTML}{0288d1}

% grey
\definecolor{c0}{HTML}{f5f5f5}
\definecolor{c1}{HTML}{e0e0e0}
\definecolor{c2}{HTML}{9e9e9e}
\definecolor{c3}{HTML}{616161}
\definecolor{c4}{HTML}{424242}

\tikzset{
    % nodegray/.style={circle,draw,fill=c1!10},
    archnode/.style={thick,circle,c4,draw=c4,fill=c0},
    % archnode/.style={circle,draw,fill=grey!20}
    % nodeblue/.style={circle,draw,fill=c3},
}
\begin{tikzpicture}[
outer/.style={draw=gray,dashed,fill=green!1,thick},
>={Stealth[inset=0pt,
length=4pt,angle'=45,round]},scale=1.,
sty edge sel/.style={ultra thick,color=c4},
sty node sel/.style={ultra thick,circle,c4,draw=c4,fill=c0,minimum size=\sizenode},
sty node unsel/.style={circle,c2,draw=c2,fill=c0,minimum size=\sizenode},
sty edge unsel/.style={color=c2},
]

\def\sizenode{10};

\node[sty node sel] (0) at (1,1) {0};
\node[sty node sel] (1) at (0,0) {1};
\node[sty node sel] (2) at (1,-1) {2};
\node[sty node sel] (3) at (2,0.2) {3};
\node[sty node unsel] (4) at (3.0,0.9) {4};
\node[sty node sel] (5) at (2.8,-0.7) {5};
\node[sty node unsel] (6) at (3.9,0.1) {6};
\node[sty node unsel] (8) at (4.9,0.8) {8};
\node[sty node unsel] (9) at (5.6,-0.2) {9};
\node[sty node sel] (7) at (4.5,-1) {7};

\def\gcnx{8.25};
\def\gcny{-1.1};
\def\offsetx{1};
\def\offsety{1};
\def\sizenodegcn{6};
\def\scalegcn{0.9};

% layer 1
\node[archnode,minimum size=\sizenodegcn,scale=\scalegcn] (g00) at (\gcnx,\gcny) {0};
\node[archnode,minimum size=\sizenodegcn,scale=\scalegcn] (g01) at (\gcnx+\offsetx,\gcny) {1};
\node[archnode,minimum size=\sizenodegcn,scale=\scalegcn] (g02) at (\gcnx+2*\offsetx,\gcny) {2};
\node[archnode,minimum size=\sizenodegcn,scale=\scalegcn] (g03) at (\gcnx+3*\offsetx,\gcny) {3};
\node[archnode,minimum size=\sizenodegcn,scale=\scalegcn] (g05) at (\gcnx+4*\offsetx,\gcny) {5};
\node[archnode,minimum size=\sizenodegcn,scale=\scalegcn] (g07) at (\gcnx+5*\offsetx,\gcny) {7};

% layer 2
\node[archnode,minimum size=\sizenodegcn,scale=\scalegcn] (g10) at (\gcnx,\gcny+\offsety) {0};
\node[archnode,minimum size=\sizenodegcn,scale=\scalegcn] (g11) at (\gcnx+\offsetx,\gcny+\offsety) {1};
\node[archnode,minimum size=\sizenodegcn,scale=\scalegcn] (g12) at (\gcnx+2*\offsetx,\gcny+\offsety) {2};
\node[archnode,minimum size=\sizenodegcn,scale=\scalegcn] (g13) at (\gcnx+3*\offsetx,\gcny+\offsety) {3};
\node[archnode,minimum size=\sizenodegcn,scale=\scalegcn] (g15) at (\gcnx+4*\offsetx,\gcny+\offsety) {5};
\node[archnode,minimum size=\sizenodegcn,scale=\scalegcn] (g17) at (\gcnx+5*\offsetx,\gcny+\offsety) {7};

% layer 3
\node[archnode,minimum size=\sizenodegcn,scale=\scalegcn] (g20) at (\gcnx,\gcny+2*\offsety) {0};
\node[archnode,minimum size=\sizenodegcn,scale=\scalegcn] (g21) at (\gcnx+\offsetx,\gcny+2*\offsety) {1};
\node[archnode,minimum size=\sizenodegcn,scale=\scalegcn] (g22) at (\gcnx+2*\offsetx,\gcny+2*\offsety) {2};
\node[archnode,minimum size=\sizenodegcn,scale=\scalegcn] (g23) at (\gcnx+3*\offsetx,\gcny+2*\offsety) {3};
\node[archnode,minimum size=\sizenodegcn,scale=\scalegcn] (g25) at (\gcnx+4*\offsetx,\gcny+2*\offsety) {5};
\node[archnode,minimum size=\sizenodegcn,scale=\scalegcn] (g27) at (\gcnx+5*\offsetx,\gcny+2*\offsety) {7};

% 1-2 link
\draw [-latex,c4,fill=c4] (g00.north) -- (g10.south);
\draw [-latex,c4,fill=c4] (g00.north) -- (g11.south);
\draw [-latex,c4,fill=c4] (g00.north) -- (g13.south);
\draw [-latex,c4,fill=c4] (g01.north) -- (g11.south);
\draw [-latex,c4,fill=c4] (g01.north) -- (g10.south);
\draw [-latex,c4,fill=c4] (g01.north) -- (g12.south);
\draw [-latex,c4,fill=c4] (g01.north) -- (g13.south);
\draw [-latex,c4,fill=c4] (g02.north) -- (g11.south);
\draw [-latex,c4,fill=c4] (g02.north) -- (g12.south);
\draw [-latex,c4,fill=c4] (g02.north) -- (g13.south);
\draw [-latex,c4,fill=c4] (g03.north) -- (g10.south);
\draw [-latex,c4,fill=c4] (g03.north) -- (g11.south);
\draw [-latex,c4,fill=c4] (g03.north) -- (g12.south);
\draw [-latex,c4,fill=c4] (g03.north) -- (g13.south);
\draw [-latex,c4,fill=c4] (g03.north) -- (g15.south);
\draw [-latex,c4,fill=c4] (g05.north) -- (g13.south);
\draw [-latex,c4,fill=c4] (g05.north) -- (g15.south);
\draw [-latex,c4,fill=c4] (g05.north) -- (g17.south);
\draw [-latex,c4,fill=c4] (g07.north) -- (g15.south);
\draw [-latex,c4,fill=c4] (g07.north) -- (g17.south);

% 2-3 link
\draw [-latex,c4,fill=c4] (g10.north) -- (g20.south);
\draw [-latex,c4,fill=c4] (g10.north) -- (g21.south);
\draw [-latex,c4,fill=c4] (g10.north) -- (g23.south);
\draw [-latex,c4,fill=c4] (g11.north) -- (g21.south);
\draw [-latex,c4,fill=c4] (g11.north) -- (g20.south);
\draw [-latex,c4,fill=c4] (g11.north) -- (g22.south);
\draw [-latex,c4,fill=c4] (g11.north) -- (g23.south);
\draw [-latex,c4,fill=c4] (g12.north) -- (g21.south);
\draw [-latex,c4,fill=c4] (g12.north) -- (g22.south);
\draw [-latex,c4,fill=c4] (g12.north) -- (g23.south);
\draw [-latex,c4,fill=c4] (g13.north) -- (g20.south);
\draw [-latex,c4,fill=c4] (g13.north) -- (g21.south);
\draw [-latex,c4,fill=c4] (g13.north) -- (g22.south);
\draw [-latex,c4,fill=c4] (g13.north) -- (g23.south);
\draw [-latex,c4,fill=c4] (g13.north) -- (g25.south);
\draw [-latex,c4,fill=c4] (g15.north) -- (g23.south);
\draw [-latex,c4,fill=c4] (g15.north) -- (g25.south);
\draw [-latex,c4,fill=c4] (g15.north) -- (g27.south);
\draw [-latex,c4,fill=c4] (g17.north) -- (g25.south);
\draw [-latex,c4,fill=c4] (g17.north) -- (g27.south);


\path
    (1) [bend left] edge [sty edge sel] node [right] {} (0)
    (0) [bend left] edge [sty edge sel] node [right] {} (3)
    (1) [bend right] edge [sty edge sel] node [right] {} (2)
    (2) [bend right] edge [sty edge sel] node [right] {} (3)
    (1) [bend right=10] edge [sty edge sel] node [right] {} (3)
    (3) [bend left] edge [sty edge unsel] node [right] {} (4)
    (4) [bend left=20] edge [sty edge unsel] node [right] {} (8)
    (8) [bend left] edge [sty edge unsel] node [right] {} (9)
    (7) [bend right] edge [sty edge unsel] node [right] {} (9)
    (3) [bend right] edge [sty edge sel] node [right] {} (5)
    (5) [bend right] edge [sty edge sel] node [right] {} (7)
    (4) [bend right] edge [sty edge unsel] node [right] {} (6)
    (6) [bend right] edge [sty edge unsel] node [right] {} (8)
    (8) [bend left=15] edge [sty edge unsel] node [right] {} (7);

% text and arrow
\node[] at (2.8,-1.9) {$\G[s]=\mbox{\sample}(\G)$};
\node[] at (\gcnx+2.5*\offsetx,-1.9) {Full GCN on $\G[s]$};
% \draw[->,>=latex,blue!20!white,line width=3pt,bend left=70] (6,0.3) to (7.7,0.3);
% \draw[->,>=latex,blue!20!white,line width=3pt,bend left=70] (7.7,-0.3) to (6,-0.3);
\draw[-to,line width=2pt,c4] (6.9,0.1)+(170:0.75) arc(170:20:0.75);
\draw[-to,line width=2pt,c4] (6.9,-0.1)+(-10:0.75) arc(-10:-160:0.75);

\end{tikzpicture}

\caption{{\graphsaint} training algorithm}
    \label{fig: gsaint overall}
\end{figure}
%[1.25 pages]


\begin{algorithm}
\caption{{\graphsaint} training algorithm}
\label{algo: gsaint meta}
\begin{algorithmic}[1]
%\label{algo: gsaint meta}
\renewcommand{\algorithmicrequire}{\textbf{Input:}}
\renewcommand{\algorithmicensure}{\textbf{Output:}}
\Require Training graph $\G\paren{\V,\E,\X}$; Labels $\overline{\bm{Y}}$; Sampler \sample;
\Ensure GCN model with trained weights
%\For{each phase}
\State Pre-processing: Setup {\sample} parameters; Compute normalization coefficients $\alpha$, $\lambda$.
\For{each minibatch}
    \State $\G[s]\paren{\V[s],\E[s]} \gets$ Sampled sub-graph of $\G$ according to {\sample}
    \State GCN construction on $\G[s]$.
    \State $\set*{\bm{y}_v\given v\in\V[s]}\gets$ Forward propagation of $\set*{\bm{x}_v\given v\in \V[s]}$, normalized by $\alpha$
    %\State Loss calculation on $\bm{y}_v$ and $\overline{\bm{y}}_v$, $v\in\V[s]$.
    \State Backward propagation from $\lambda$-normalized loss $L\paren{\bm{y}_v,\overline{\bm{y}}_v}$. Update weights. 
\EndFor
%\EndFor
%\State \Return $\Set{\setwl{i}|1\leq i \leq D}$
\end{algorithmic} 
\end{algorithm}


{\graphsaint} follows the design philosophy of directly sampling the training graph $\G$, rather than the corresponding GCN. 
Our goals are to 
\begin{enumerate*}
    \item extract appropriately connected subgraphs so that little information is lost when propagating within the subgraphs, and
    \item combine information of many subgraphs together so that the training process overall learns good representation of the full graph.
\end{enumerate*}
% Intuitively, using proper graph sampling algorithms, we can extract small representative subgraphs from the large training graph, such that accurate node embedding can be obtained by propagating information \textit{within the subgraphs}. 
%The major benefit of graph sampling based training is, it allows much flexibility in GCN architecture design (see Section \ref{sec: ext}). In short, 


Figure \ref{fig: gsaint overall} and Algorithm \ref{algo: gsaint meta} illustrate the training algorithm. 
Before training starts, we perform light-weight pre-processing on $\G$ with the given sampler \sample. The pre-processing estimates the probability of a node $v\in\V$ and an edge $e\in\E$ being sampled by {\sample}. Such probability is later used to normalize  the subgraph neighbor aggregation and the minibatch loss (Section \ref{sec: bias}). Afterwards, training proceeds by iterative weight updates via SGD. 
Each iteration starts with an independently sampled $\G[s]$ (where $\size{\V[s]}\ll \size{\V}$). We then build a full GCN on $\G[s]$ to generate embedding and calculate loss for every $v\in \V[s]$. 
In Algorithm \ref{algo: gsaint meta}, node representation is learned by performing node classification in the supervised setting, and each training node $v$ comes with a ground truth label $\overline{\bm{y}}_v$. 

%One would naturally wonder how to design the function {\sample}. 
Intuitively, there are two requirements for {\sample}: 
\begin{enumerate*}
    \item Nodes having high influence on each other should be sampled in the same subgraph.
    \item Each edge should have non-negligible probability to be sampled. 
\end{enumerate*}
For requirement (1), an ideal {\sample} would consider the joint information from node connections as well as attributes. However, the resulting algorithm may have high complexity as it would need to infer the relationships between features. For simplicity, we define ``influence'' from the graph connectivity perspective and design topology based samplers. 
Requirement (2) leads to better generalization since it enables the neural net to explore the full feature and label space. % with low sample complexity...
%Note that {\sample} satisfying requirement (1) almost inevitably introduces bias towards ``high centrality'' nodes. Such bias can be addressed by properly normalizing the minibatch loss, as shown in Section \ref{sec: method-loss}.


%The derivation on ``no bias'' assumes at least one neighbor for each subgraph node. 



\subsection{Normalization}
\label{sec: bias}

%%Aggregation normalization $\alpha$ (edge probability conditioned on node probability); loss normalization $\lambda$ (node probability) -- these lead to an unbiased estimator.

%Normalization of $\alpha$ (aggregator normalization), $\lambda$ (loss normalization) to guarantee unbiased estimator. Also describe the pre-processing step. 

A sampler that preserves connectivity characteristic of $\G$ will almost inevitably introduce bias into minibatch estimation. 
In the following, we present normalization techniques to eliminate biases. 

Analysis of the complete multi-layer GCN is difficult due to non-linear activations. Thus, we analyze the embedding of each layer independently. 
This is similar to the treatment of layers independently by prior work \citep{fastgcn,as-gcn}.
Consider a layer-$\paren{\ell+1}$ node $v$ and a layer-$\ell$ node $u$. If $v$ is sampled (i.e., $v\in\V[s]$), we can compute the aggregated feature of $v$ as:
\begin{equation}
    \zeta_{v}^{(\ell+1)} = \sum\limits_{u\in\V}\frac{\Anorm_{v,u}}{\alpha_{u,v}}
    \paren{\W[\ell]}^{\trans} \bm{x}_u^{\paren{\ell}}\mathbbm{1}_{u|v} = 
    \sum\limits_{u\in\V}\frac{\Anorm_{v,u}}{\alpha_{u,v}} \Tilde{\bm{x}}_u^{\paren{\ell}} \mathbbm{1}_{u|v},
\end{equation}
where $\Tilde{\bm{x}}_u^{\paren{\ell}} = \paren{\W[\ell]}^{\trans}  \bm{x}_u^{\paren{\ell}}$, and $\mathbbm{1}_{u |v}\in\set{0,1}$ is the indicator function given $v$ is in the subgraph (i.e., $\mathbbm{1}_{u|v}=0$ if $v\in\V[s]\wedge \paren{u,v}\not\in\E[s]$; $\mathbbm{1}_{u|v}=1$ if $\paren{u,v}\in\E[s]$; $\mathbbm{1}_{u|v}$ not defined if $v\not\in\V[s]$). We refer to the constant $\alpha_{u,v}$ as \emph{aggregator normalization}.
Define $p_{u,v}=p_{v,u}$ as the probability of an edge $\paren{u,v}\in\E$ being sampled in a subgraph, and $p_v$ as the probability of a node $v\in\V$ being sampled.
\begin{proposition}
\label{prop: norm}
$\zeta_{v}^{(\ell+1)}$ is an unbiased estimator of the aggregation of $v$ in the full $\paren{\ell+1}^{\text{th}}$ GCN layer, if $\alpha_{u,v} = \frac{p_{u,v}}{p_v}$. i.e.,  
$
    \mathbb{E}\paren{\zeta_{v}^{(\ell+1)}} = \sum\limits_{u\in\V}\Anorm_{v,u} \Tilde{\bm{x}}_u^{\paren{\ell}}\nonumber
$.
\end{proposition}

%We use Proposition~\ref{prop: norm} considering each layer independently for aggregator normalization during feature propagation in \graphsaint, assuming that each layer independently learns an embedding.
Assuming that each layer independently learns an embedding, we use Proposition~\ref{prop: norm} to normalize feature propagation of each layer of the GCN built by \graphsaint.
Further, let $L_v$ be the loss on $v$ in the output layer. The minibatch loss is calculated as $L_{\text{batch}} = \sum_{v \in \G[s]} L_v/\lambda_v$, where $\lambda_v$ is a constant that we term \emph{loss normalization}. We set $\lambda_v = \size{\V}\cdot p_v$ so that:
\begin{equation}
    \mathbb{E}\paren{L_{\text{batch}}} = \frac{1}{\size{\mathbb{G}}}\sum_{\G[s]\in\mathbb{G}}\sum_{v\in\V[s]} \frac{L_v}{\lambda_v} = \frac{1}{\size{\V}}\sum_{v\in\V} L_v.
\end{equation}

Feature propagation within subgraphs thus requires normalization factors $\alpha$ and $\lambda$, which are computed based on the edge and node probability $p_{u,v}$, $p_v$. 
In the case of random node or random edge samplers, $p_{u,v}$ and $p_v$ can be derived analytically. For other samplers in general, closed form expression is hard to obtain. Thus, we perform pre-processing for estimation. 
Before training starts, we run the sampler repeatedly to obtain a set of $N$ subgraphs $\mathbb{G}$. We setup a counter $C_v$ and $C_{u,v}$ for each $v\in\V$ and $\paren{u,v}\in\E$, to count the number of times the node or edge appears in the subgraphs of $\mathbb{G}$. 
Then we set $\alpha_{u,v}=\frac{C_{u,v}}{C_v}=\frac{C_{v,u}}{C_v}$ and $\lambda_v=\frac{C_v}{N}$.
The subgraphs $\G[s]\in\mathbb{G}$ can all be reused as minibatches during training. Thus, the overhead of pre-processing is small (see Appendix \ref{appendix: sampler cost}). 
%The space complexity is 

%%%%Therefore, with aggregate normalization $\alpha$ and loss normalization $\lambda$, minibatch loss of {\graphsaint} $L_{batch}$ is an unbiased estimator of the full batch loss. 


% For random walk sampler, let $m$ be the number of independent walkers, $k$ be the depth of each walker, and $\bm{B}$ be the transition matrix (i.e., $\bm{B}=\bm{D}^{-1}\bm{A}=\Anorm$). Let $\bm{D}$ be the diagonal degree matrix (i.e., $\bm{D}_{u,u}=\text{deg}\paren{u}$), and $\bm{I}$ be an identity vector. 

% \begin{subequations}
% \begin{align}
%     \bm{r} =& \frac{1}{\size{\V}}\bm{D}^{-1}\paren{\bm{B}^0+\bm{B}^1+\cdots+\bm{B}^{k-1}}\bm{I}\\
%     P_u =& 1 - \paren{1-\bm{r}_u}^m\\
%     P_{e} =& 1-\paren{1-\paren{\bm{r}_u+\bm{r}_v}}^m
% \end{align}
% \end{subequations}

%So $\lambda_u=P_u$ and $\alpha_{e|u}=P_e/P_u$. 

\subsection{Variance}
\label{sec: var}
We derive samplers for variance reduction. Let $e$ be the edge connecting $u$, $v$, and $\bm{b}_e^{(\ell)} = \Anorm_{v, u}\Tilde{\bm{x}}_u^{(\ell-1)} + \Anorm_{u, v}\Tilde{\bm{x}}_v^{(\ell-1)}$. It is desirable that variance of all estimators $\zeta_{v}^{(\ell)}$ is small. With this objective, we define:
\begin{equation}
    \zeta =  \sum_{\ell} \sum_{v \in \G[s]} \frac{\zeta_{v}^{(\ell)}}{p_v} = \sum_{\ell} \sum_{v, u}\frac{\Anorm_{v,u}}{p_v\alpha_{u,v}} \Tilde{\bm{x}}_u^{\paren{\ell}} \mathbbm{1}_{v}\mathbbm{1}_{u|v} = \sum_{\ell}\sum_e \frac{\bm{b}_e^{(\ell)}}{p_e} \mathbbm{1}_{e}^{\paren{\ell}}.
\end{equation}
where $\mathbbm{1}_{e}=1$ if $e\in\E[s]$; $\mathbbm{1}_{e}=0$ if $e\not\in\E[s]$. And $\mathbbm{1}_v=1$ if $v\in\V[s]$; $\mathbbm{1}_v=0$ if $v\not\in\V[s]$. The factor $p_u$ in the first equality is present so that $\zeta$ is an unbiased estimator of the sum of all node aggregations at all layers: $\mathbb{E}\paren{\zeta} = \sum_\ell \sum_{v\in\V} \mathbb{E}\paren{\zeta_v^{(\ell)}}$. Note that $\mathbbm{1}_{e}^{\paren{\ell}} = \mathbbm{1}_e, \forall \ell$, since once an edge is present in the sampled graph, it is present in all layers of our GCN. 

We define the optimal edge sampler to minimize variance for every dimension of $\zeta$. We restrict ourselves to independent edge sampling. For each $e\in\E$, we make independent decision on whether it should be in $\G[s]$ or not. The probability of including $e$ is $p_e$. We further constrain $\sum p_e=m$, so that the expected number of sampled edges equals to $m$. The budget $m$ is a given sampling parameter. 

%Comparison with the layer-sampling counterpart -- show variance reduction.

%Insights -- graph sampling reduces variance due to better connectivity (less stochasticity in sampling, less dropped nodes). 
\begin{theorem}
\label{thm: var}
Under independent edge sampling with budget $m$, the optimal edge probabilities to minimize the sum of variance of each $\zeta$'s dimension  is given by: $p_e=\frac{m}{\sum_{e'}\left\lVert\sum_{\ell}\bm{b}_{e'}^{\paren{\ell}}\right\rVert}\left\lVert\sum_{\ell}\bm{b}_e^{\paren{\ell}}\right\rVert$.
\end{theorem}

To prove Theorem \ref{thm: var}, we make use of the independence among graph edges, and the dependence among layer edges to obtain the covariance of $\mathbbm{1}_e^{\paren{\ell}}$. Then using the fact that sum of $p_e$ is a constant, we use the Cauchy-Schwarz inequality to derive the optimal $p_e$. Details are in Appendix \ref{appendix: proof}. 

% In independent edge sampling, $Cov\paren{I_{e_1}^{\paren{\ell_1}}, I_{e_2}^{\paren{\ell_2}}} = 0, \forall e_1 \neq e_2$, and $Cov\paren{I_{e}^{\paren{\ell_1}}, I_{e}^{\paren{\ell_2}}} = p_e - p_e^2$. Therefore,
% \begin{align}
%     Var(\zeta) &= \sum_{e, \ell} \left(\frac{b_e^{(\ell)}}{p_e}\right)^2 Var\paren{I_{e}^{\paren{\ell}}} + 2\sum_{e, \ell_1, \ell_2} \frac{b_{e}^{\paren{\ell_1}}b_{e}^{\paren{\ell_2}}}{p_e^2}Cov\paren{I_{e}^{\paren{\ell_1}}, I_{e}^{\paren{\ell_2}}} \nonumber \\
%     &= \sum_{e, \ell} \frac{\paren{b_{e}^{\paren{\ell}}}^2}{p_e} - \sum_e \paren{b_{e}^{\paren{\ell}}}^2 + 2 \sum_{e, \ell_1, \ell_2} \frac{b_{e}^{\paren{\ell_1}}b_{e}^{\paren{\ell_2}}}{p_e^2}(p_e - p_e^2) \nonumber \\
%     &= \sum_e\frac{\paren{\sum_\ell b_e^{(\ell)}}^2}{p_e} - \sum_e \left(\sum_\ell b_e^{(\ell)}\right)^2
% \end{align}
% Let $m = \sum_e p_e$ be the expected number of edges in the sampled subgraph.
% By Cauchy-Schwarz inequality $\sum_e \left(\frac{\sum_\ell b_e^{(\ell)}}{\sqrt{p_e}}\right)^2 \sum_e \paren{\sqrt{p_e}}^2\geq \left(\sum_{e, \ell} b_e^{(\ell)}\right)^2$. 
% The equality is achieved when $b_e^{(\ell)}/\sqrt{p_e} \propto \sqrt{p_e}$, i.e., variance is minimized when
% $p_e \propto \sum_\ell b_e^{(\ell)}$.

Note that calculating $\bm{b}_e^{(\ell)}$ requires computing $\Tilde{\bm{x}}_v^{\paren{\ell - 1}}$, which increases the complexity of sampling. As a reasonable simplification, we ignore $\Tilde{\bm{x}}_v^{\paren{\ell}}$ to make the edge probability dependent on the graph topology only. Therefore, we choose $p_e \propto \Anorm_{v,u} + \Anorm_{u, v}=\frac{1}{\text{deg}\paren{u}}+\frac{1}{\text{deg}\paren{v}}$.

The derived optimal edge sampler agrees with the intuition in Section \ref{sec: gsaint}. If two nodes $u$, $v$ are connected and they have few neighbors, then $u$ and $v$ are likely to be influential to each other. In this case, the edge probability $p_{u,v}=p_{v,u}$ should be high. The above analysis on edge samplers also inspires us to design other samplers, which are presented in Section \ref{sec: sampler}.


\paragraph{Remark} We can also apply the above edge sampler to perform layer sampling. Under the independent layer sampling assumption of \cite{fastgcn}, one would sample a connection $\paren{u^{\paren{\ell}},v^{\paren{\ell+1}}}$ with probability $p_{u,v}^{\paren{\ell}}\propto \frac{1}{\text{deg}\paren{u}}+\frac{1}{\text{deg}\paren{v}}$. 
For simplicity, assume a uniform degree graph (of degree $d$). Then $p_e^{\paren{\ell}}=p$. For an already sampled $u^{\paren{\ell}}$ to connect to layer $\ell+1$, at least one of its edges has to be selected by the layer $\ell+1$ sampler. 
Clearly, the probability of an input layer node to ``survive'' the $L$ number of independent sampling process is $\paren{1-\paren{1-p}^d}^{L-1}$. 
Such layer sampler potentially returns an overly sparse minibatch for $L>1$. On the other hand, connectivity within a minibatch of {\graphsaint} never drops with GCN depth. If an edge is present in layer $\ell$, it is present in all layers. 

\subsection{Samplers}
\label{sec: sampler}

Based on the above variance analysis, we present several light-weight and efficient samplers that {\graphsaint} has integrated. Detailed sampling algorithms are listed in Appendix \ref{sec: sampler algo}.

\paragraph{Random node sampler} 
We sample $\size{\V[s]}$ nodes from $\V$ randomly, according to a node probability distribution $P\paren{u}\propto \norm{\Anorm_{:,u}}^2$. This sampler is inspired by the layer sampler of \cite{fastgcn}.


\paragraph{Random edge sampler} We perform edge sampling as described in Section \ref{sec: var}. 


\paragraph{Random walk based samplers} 
Another way to analyze graph sampling based multi-layer GCN is to ignore activations. In such case, $L$ layers can be represented as a single layer with edge weights given by $\bm{B} = \Anorm^L$. Following a similar approach as Section~\ref{sec: var}, if it were possible to pick pairs of nodes (whether or not they are directly connected in the original $\Anorm$) independently, then we would set $p_{u, v} \propto \bm{B}_{u,v} + \bm{B}_{v,u}$, where $\bm{B}_{u, v}$ can be interpreted as the probability of a random walk to start at $u$ and end at $v$ in $L$ hops (and $\bm{B}_{v, u}$ vice-versa). Even though it is not possible to sample a subgraph where such pairs of nodes are independently selected, we still consider a random walk sampler with walk length $L$ as a good candidate  for $L$-layer GCNs. There are numerous random walk based samplers proposed in the literature \citep{frontier,forest-fire,sampler_survey, sampler_icde}. 
In the experiments, we implement a regular random walk sampler (with $r$ root nodes selected uniformly at random and each walker goes $h$ hops), and also a multi-dimensional random walk sampler defined in \cite{frontier}.

% Intuitively, given the interpretation of variance, if the sample size $m\ll M$, covariance of $X_i$, $X_j$ is negative?
% Consider the special case of length-2 random walk on uniform degree graph. You can verify the value of covariance is negative. 
% Expanding from that, if we have a biased random walk ensuring $X_e\propto b_e$, then the self-variance term is the same compared with the optimal independent edge sampler. The covariance term is negative. We can then prove such random walk sampler is better than the optimal independent edge sampler. 
% Consider the uniform degree graph, such biased random walk becomes uniform random walk, given that stationary distribution of uniform random walk approaches node degree. 
% Examples: forest fire \cite{forest-fire}, Frontier \cite{frontier}



For all the above samplers, we return the subgraph induced from the sampled nodes. The induction step adds more connections into the subgraph, and empirically helps improve convergence. 

%Justify the induction step...
